\begin{table}[h!]
\centering
\caption{The five highest-scoring regions after each of the $T=3$ sessions, with the utilities computed from the FTRL-DOR parameter estimates of the corresponding session, in the case of linear marginal utility functions. The utilities are expressed as percentages of the utility of the ideal alternative, which scores 100 on every criterion. The reference alternatives are written in bold.}
\label{tab:rank_lin_top5}
\begin{tabularx}{\textwidth}{c|Cc|Cc|Cc}
\toprule
\multirow{2}{*}{Rank}&\multicolumn{2}{c}{$T=1$}&\multicolumn{2}{c}{$T=2$}&\multicolumn{2}{c}{$T=3$}\\
\cmidrule(lr){2-3}\cmidrule(lr){4-5}\cmidrule(lr){6-7}
&Region&$U\left(a\right)$&Region&$U\left(a\right)$&Region&$U\left(a\right)$\\
\midrule
1 & Veneto & 96.09 & Veneto & 96.46 & Veneto & 96.94 \\
2 & Tuscany & 95.31 & Tuscany & 95.22 & \textbf{Tuscany} & \textbf{95.12} \\
3 & Emilia Romagna & 92.80 & \textbf{Emilia Romagna} & \textbf{93.22} & \textbf{Emilia Romagna} & \textbf{94.10} \\
4 & A. P. of Trento & 92.74 & A. P. of Trento & 92.72 & A. P. of Trento & 93.58 \\
5 & Piedmont & 90.13 & Piedmont & 90.69 & Piedmont & 91.41 \\
\bottomrule
\end{tabularx}
\end{table}
